\documentclass[11pt,a4paper]{article}

\usepackage[margin=2.5cm]{geometry}
\usepackage{hyperref}
\usepackage{amsmath,amssymb}
\usepackage{parskip}
\usepackage{microtype}
\usepackage{titlesec}
\usepackage{booktabs}
\usepackage{array}
\usepackage{enumitem}
\usepackage{graphicx}

\titleformat{\section}{\large\bfseries}{\thesection}{0.6em}{}
\titleformat{\subsection}{\normalsize\bfseries}{\thesubsection}{0.6em}{}

\newcommand{\smax}{\sigma_{\max}}
% ---- Provenance (comments only; nothing of this appears in the PDF) ----------------
% Every table is \input from a fragment generated by scripts/train_render.py (the
% training record and the timing run) or scripts/frozen_render.py (the checkpoints on
% the held-out tasks); each fragment's header names the result files, the code commit
% and the Slurm job. Numbers in the text are read off those tables.
% The training run: code commit 654ad1bc641fbd42a46f392fd335fde2e8c134f6, Slurm job
%   3460844 (train-2replies), one A30 on gx17v1, 2026-11-25 22:44 to 2026-11-26 11:08,
%   config configs/loop_grpo_branched_crystallizer.yaml, default seed, submitted with
%   --export=ALL,LORA4VR_GIT_COMMIT=<commit> from its own code directory on the cluster
%   (~/work/lora4vr-train-654ad1b). Record: docs/training/train-2replies_3460844/train/
%   (the config, progress/trainer_log_iter_NNN.json with one entry per step, the log
%   without the tables of completions, and completions.tgz, every conversation of the
%   run as the trainer saved it). Tables: docs/training/train-2replies_3460844/tables/.
%   The 23 adapters are on the cluster only:
%   ~/work/lora4vr-train-654ad1b/models/lora/train-2replies_3460844/iter_000 ... iter_022.
% The timing run: code commit 50df6169756a9047299450d1b305c54a470715a1, Slurm job
%   3460791 (train-timing-a30), config configs/timing_grpo_branched_crystallizer.yaml.
%   Results: docs/train_timing/train-timing-a30_3460791/.
% The checkpoints on the held-out tasks: code commit
%   b80823043e95e55a41fd518734c287892743d7c4; Slurm jobs 3461865 (checkpoint 1), 3462056
%   (checkpoint 2), 3462280 (checkpoint 3), one A30 each, submitted with
%   --export=ALL,LORA4VR_GIT_COMMIT=<commit>,N_ADAPTERS=<k>,LORA_DIR=<the adapter directory>.
%   Results: docs/training/train-2replies_3460844/eval-train-2replies-ckpt<kk>_<job>/
%   (job log, stage2_summary.json with its provenance block and the hash of every
%   adapter merged, stage2_samples.jsonl with every sample). Tables:
%   docs/training/train-2replies_3460844/tables_eval/, rendered beside the frozen
%   model's Stage 2 run, job 3458956 (commit 7800206, docs/stage12_agent/).
%   Jobs 3463506 (checkpoint 5) and 3463507 (checkpoint 23) were submitted on
%   2026-11-26 and wait for the A30 node, which is in maintenance until 2026-11-28.
% ------------------------------------------------------------------------------------
% GENERATED by scripts/train_render.py. Do not edit; re-run the renderer.
% run: docs/training/train-2replies_3460844/train, code commit 654ad1bc641fbd42a46f392fd335fde2e8c134f6, Slurm job 3460844, config loop_grpo_branched_crystallizer.yaml, 23 iterations, 552 steps
% timing run: docs/train_timing/train-timing-a30_3460791, code commit 50df6169756a9047299450d1b305c54a470715a1, Slurm job 3460791, 4 steps
\newcommand{\TRcommit}{654ad1b}
\newcommand{\TRjob}{3460844}
\newcommand{\TRiterations}{23}
\newcommand{\TRsteps}{552}
\newcommand{\TRstepsperiteration}{24}
\newcommand{\TRgone}{4}
\newcommand{\TRgtwo}{4}
\newcommand{\TRconversations}{16}
\newcommand{\TRmaxturntokens}{1000}
\newcommand{\TRrank}{16}
\newcommand{\TRbestiteration}{5}
\newcommand{\TRbestreward}{$-$26.6}
\newcommand{\TRlastreward}{$-$54.3}
\newcommand{\TRhours}{12.4}
\newcommand{\TRtimingcommit}{50df616}
\newcommand{\TRtimingjob}{3460791}


\title{\textbf{The first training run: two replies, reward only}}
\author{}
\date{2026-11-26}

\begin{document}
\maketitle

\noindent
The frozen 4B model, given two replies and the tool, calls the probe run in
nine samples of ten and then, in seven of ten, probes again with its last
reply and never hands in a controller (\texttt{docs/stage12\_agent\_loop.tex}).
Stating the replies left did not change that. The question for training was
therefore open in both directions: reward alone might teach the model to
answer after its probe, or it might teach it that the probe is not worth a
reply.

This document records the first training run on that setup: two replies,
the reward as the only signal, the brief untouched. In its first
\TRstepsperiteration{} gradient steps the model learned to answer after the
probe. In the next seventy it dropped the probe. For the remaining four
fifths of the run it trained a controller that does not know the hidden
limit, and did not get that controller near the cautious hand-written one.
The run is one seed and settles nothing about the method in general; it
shows what the reward says at this starting point, and why.

\paragraph{Terms.} A \emph{conversation} is one task played once: the brief,
the model's first \emph{reply}, and, if that reply was not the final answer,
the tool response and a second reply. The \emph{probe run} is the tool
\texttt{probe\_run}; the \emph{final answer} is a reply without a tool call,
and the \emph{final run} is its controller on the scored batch. A
\emph{step} is one gradient update, made from the conversations of one task.
An \emph{iteration} is one pass over the 24 training tasks, so
\TRstepsperiteration{} steps. An \emph{adapter} is the small set of trained
weights (LoRA, rank~\TRrank) an iteration saves. \emph{Checkpoint $k$} is the
model after $k$ iterations: the base model with the first $k$ adapters
merged in; checkpoint~0 is the frozen model. The \emph{error floor} is the
reward of a conversation that ends without a final run or whose controller
does not run, $-300$. \emph{Held-out tasks} are the 48 evaluation tasks,
never trained on. \emph{Success}, \emph{fast success} (success within
$1.25\times$ the oracle's steps), \emph{ruined} (a shower or lost seeds) and
the \emph{ride ratio} (the highest reported $\sigma$ of the final run over
the batch's limit) are as in the Stage~1 and~2 report. The \emph{cautious
blind controller} and the \emph{oracle} are the hand-written references:
the first is safe for every limit in the range, the second knows the limit.

\section{The run}

\paragraph{Method.} Each step takes one training task and builds a tree of
conversations: \TRgone{} first replies sampled from the brief, and for each
first reply that calls the probe run, \TRgtwo{} second replies sampled after
its tool response, so \TRconversations{} conversations per step. A first
reply that is already the final answer is a leaf: it has one reward, counted
\TRgtwo{} times. Every second reply is scored as the final answer; a second
reply that calls the tool again is at the error floor. Credit is given per
reply: a second reply is credited with its reward less the mean reward of
the \TRgtwo{} second replies that share its first reply; a first reply is
credited with the mean reward of its own second replies less the mean over
the whole tree. The loss is the clipped surrogate of GRPO on these credits.
There is no other signal: no example, no imitation, no bonus for probing.

\paragraph{Sizing.} \TRiterations{} iterations of \TRstepsperiteration{}
steps, \TRsteps{} steps in all; at most \TRmaxturntokens{} new tokens per
reply; learning rate $10^{-4}$ with a warm-up over the first tenth of each
iteration and a linear decay to its end; sampling temperature 0.8. Each
iteration trains a new adapter on the model with all earlier adapters merged
in, so the learning-rate schedule and the optimizer start anew every
\TRstepsperiteration{} steps. The config is
\texttt{configs/loop\_grpo\_branched\_crystallizer.yaml}. In the experiment
plan this is the sizing of Stage~6; Stages~3 to~5 (training with one reply,
and the two simpler ways of giving credit) have not been run.

\paragraph{Evaluation is apart from training.} The loop used to evaluate
the held-out tasks after every iteration. One such evaluation takes about as
long as an iteration, and nothing in training reads it, so the loop now has
an evaluation period (\texttt{loop.eval\_every}), set longer than the run.
Every iteration saves its adapter, and a checkpoint is evaluated by a
separate job on another GPU (\texttt{jobs/hpi/eval\_checkpoint.sbatch}):
the script of Stage~2, \texttt{scripts/eval\_frozen\_baseline.py}, with the
checkpoint's adapters merged into the base model, so that a checkpoint's
numbers stand next to the frozen model's. The job refuses a checkpoint whose
adapters are not on disk in full, and the result names the hash of every
adapter it merged.

\paragraph{Cost.} Before the run, four steps at the run's sizing were timed
on one A30 (24~GB):

% provenance: see the header of training/train-2replies_3460844/tables/timing.tex
\begin{table}[h]
\centering\small
% GENERATED by scripts/train_render.py. Do not edit; re-run the renderer.
% run: docs/training/train-2replies_3460844/train, code commit 654ad1bc641fbd42a46f392fd335fde2e8c134f6, Slurm job 3460844, config loop_grpo_branched_crystallizer.yaml, 23 iterations, 552 steps
% timing run: docs/train_timing/train-timing-a30_3460791, code commit 50df6169756a9047299450d1b305c54a470715a1, Slurm job 3460791, 4 steps
\begin{tabular}{@{}rrrrr@{}}
\toprule
Step & Conversations & Seconds, rollout and update & Peak memory (GiB) & At the error floor \\
\midrule
1 & 16 & 206 & 17.2 & 81\% \\
2 & 16 & 238 & 17.8 & 56\% \\
3 & 16 & 173 & 17.1 & 75\% \\
4 & 16 & 153 & 14.6 & 100\% \\
\midrule
\multicolumn{5}{@{}l}{training time 915 s for 4 steps: 229 s per step, logging included} \\
\bottomrule
\end{tabular}

\caption{The timing run. A step with all \TRconversations{} conversations at
two replies takes close to four minutes and fits the A30 with 6~GiB to
spare.}
\end{table}

At that rate \TRsteps{} steps would take 35 hours. The run took \TRhours{} hours:
once the model answered with its first reply, a step was one batch of
\TRgone{} replies and no second batch (Table~\ref{tab:record}, the column of
seconds).

\section{The training record}

The trainer logs every step. Two of its counters say what the model does.
Each step makes one generate call for the \TRgone{} first replies and one
more for every first reply that did not end the task; the number of calls
less one is therefore the number of first replies that were followed by a
second reply, which is the number of probe runs, plus the rare tool call
that could not be acted on. The other counter is the share of the step's
conversations at the error floor.

% provenance: see the header of training/train-2replies_3460844/tables/phases.tex
\begin{table}[h]
\centering\small
\resizebox{\linewidth}{!}{% GENERATED by scripts/train_render.py. Do not edit; re-run the renderer.
% run: docs/training/train-2replies_3460844/train, code commit 654ad1bc641fbd42a46f392fd335fde2e8c134f6, Slurm job 3460844, config loop_grpo_branched_crystallizer.yaml, 23 iterations, 552 steps
% timing run: docs/train_timing/train-timing-a30_3460791, code commit 50df6169756a9047299450d1b305c54a470715a1, Slurm job 3460791, 4 steps
\begin{tabular}{@{}lrrrrrr@{}}
\toprule
Iterations & Steps & First replies followed by a second & At the error floor & Reward & Seconds per step & Peak memory (GiB) \\
\midrule
1 & 24 & 86 of 96 & 45\% & $-$166.0 & 183 & 17.4 \\
2 & 24 & 37 of 96 & 6\% & $-$62.4 & 132 & 17.9 \\
3 & 24 & 17 of 96 & 9\% & $-$66.9 & 89 & 16.7 \\
4 & 24 & 4 of 96 & 2\% & $-$48.4 & 58 & 15.5 \\
5 & 24 & 1 of 96 & 1\% & $-$26.6 & 54 & 16.1 \\
6--10 & 120 & 4 of 480 & 0\% & $-$41.2 & 49 & 16.6 \\
11--15 & 120 & 5 of 480 & 1\% & $-$47.8 & 51 & 16.6 \\
16--23 & 192 & 5 of 768 & 0\% & $-$55.1 & 44 & 16.4 \\
\midrule
all & 552 & 159 of 2208 & 3\% & $-$54.6 & 59 & 17.9 \\
\bottomrule
\end{tabular}
}
\caption{The training record, on the training tasks: the first five
iterations one by one, then three stretches. Every iteration is in
Table~\ref{tab:full} at the end.}
\label{tab:record}
\end{table}

Three things happen in turn.

\begin{enumerate}[leftmargin=*, itemsep=0.3em]
\item \textbf{Iteration 1: the model keeps probing and learns to answer.}
  It probes in 86 of 96 first replies, as the frozen model does. The share
  of conversations at the error floor is 45\% over the iteration and 6\%
  over the next.
\item \textbf{Iterations 2 to 4: the probe is dropped.} 37, then 17, then 4
  of 96 first replies are followed by a second. Step by step, the model
  still probes in three or four of four first replies for the first seven
  steps of iteration~2, and in at most two, mostly one or none, from its
  eighth step on.
\item \textbf{Iterations 5 to 23: no probe, and no further gain.} About one
  first reply in a hundred is followed by a second. The reward is best in
  iteration~\TRbestiteration{} (\TRbestreward) and then falls, to about $-55$
  over the last eight iterations (\TRlastreward{} in the last one).
\end{enumerate}

\section{The checkpoints on the held-out tasks}

Checkpoints~1, 2 and~3 were evaluated as the frozen model was in Stage~2:
the 48 held-out tasks, 4 samples each, two replies, 1000 tokens per reply.
The frozen model's row is its Stage~2 run. An evaluation is a fresh draw at
temperature 0.8; between two draws of the same model, counts out of 192 have
differed by up to about fifteen (\texttt{docs/stage12\_agent\_loop.tex},
Section~5).

% provenance: see the header of training/train-2replies_3460844/tables_eval/comparison.tex
\begin{table}[h]
\centering\small
\resizebox{\linewidth}{!}{% GENERATED by scripts/frozen_render.py. Do not edit; re-run the renderer.
% frozen: docs/stage12_agent/stage2-4b-agent_3458956, code commit 78002066ae308d6dc6aa0ebd95f400b154fa3d49, Slurm job 3458956, model models/hf_base/iter_000, 48 tasks x 4 samples, max_tokens 1000
% ckpt1: docs/training/train-2replies_3460844/eval-train-2replies-ckpt01_3461865, code commit b80823043e95e55a41fd518734c287892743d7c4, Slurm job 3461865, model models/hf_base/iter_000 with the first 1 adapter(s) of /sc/home/nima.hamidisiboni/work/lora4vr-train-654ad1b/models/lora/train-2replies_3460844 merged in, 48 tasks x 4 samples, max_tokens 1000
% ckpt2: docs/training/train-2replies_3460844/eval-train-2replies-ckpt02_3462056, code commit b80823043e95e55a41fd518734c287892743d7c4, Slurm job 3462056, model models/hf_base/iter_000 with the first 2 adapter(s) of /sc/home/nima.hamidisiboni/work/lora4vr-train-654ad1b/models/lora/train-2replies_3460844 merged in, 48 tasks x 4 samples, max_tokens 1000
% ckpt3: docs/training/train-2replies_3460844/eval-train-2replies-ckpt03_3462280, code commit b80823043e95e55a41fd518734c287892743d7c4, Slurm job 3462280, model models/hf_base/iter_000 with the first 3 adapter(s) of /sc/home/nima.hamidisiboni/work/lora4vr-train-654ad1b/models/lora/train-2replies_3460844 merged in, 48 tasks x 4 samples, max_tokens 1000
% ckpt5: docs/training/train-2replies_3460844/eval-train-2replies-ckpt05_3470604, code commit b80823043e95e55a41fd518734c287892743d7c4, Slurm job 3470604, model models/hf_base/iter_000 with the first 5 adapter(s) of /sc/home/nima.hamidisiboni/work/lora4vr-train-654ad1b/models/lora/train-2replies_3460844 merged in, 48 tasks x 4 samples, max_tokens 1000
% ckpt23: docs/training/train-2replies_3460844/eval-train-2replies-ckpt23_3470605, code commit b80823043e95e55a41fd518734c287892743d7c4, Slurm job 3470605, model models/hf_base/iter_000 with the first 23 adapter(s) of /sc/home/nima.hamidisiboni/work/lora4vr-train-654ad1b/models/lora/train-2replies_3460844 merged in, 48 tasks x 4 samples, max_tokens 1000
\begin{tabular}{@{}llrrrrrrrr@{}}
\toprule
Run & Model & Valid & Success & Fast & Norm.\ time & Reward & Ride ratio & Probe rate & Est.\ error \\
\midrule
frozen & iter\_000 & 29.7\% & 9.4\% & 1.0\% & 0.049 & $-$231.7 & 0.65 & 89.1\% & 0.0293 \\
ckpt1 & checkpoint 1 & 84.9\% & 22.4\% & 1.0\% & 0.118 & $-$99.2 & 0.67 & 84.9\% & 0.0071 \\
ckpt2 & checkpoint 2 & 92.2\% & 25.0\% & 2.6\% & 0.154 & $-$76.1 & 0.82 & 10.9\% & 0.0562 \\
ckpt3 & checkpoint 3 & 91.7\% & 22.4\% & 2.6\% & 0.148 & $-$80.9 & 0.89 & 27.1\% & 0.0719 \\
ckpt5 & checkpoint 5 & 98.4\% & 32.8\% & 5.2\% & 0.214 & $-$41.9 & 0.91 & 0.5\% & 0.0018 \\
ckpt23 & checkpoint 23 & 100.0\% & 25.5\% & 25.0\% & 0.233 & $-$41.1 & 1.00 & 0.0\% & -- \\
\midrule
cautious blind reference & -- & -- & -- & 4.2\% & 0.493 & 38.4 & -- & -- & -- \\
oracle & -- & -- & -- & 100.0\% & 1.000 & 79.7 & -- & -- & -- \\
\bottomrule
\end{tabular}
}
\caption{The frozen model and the first three checkpoints. Valid: the task
ended with a final run whose controller passed validation. Probe rate: the
first reply was a probe run whose output came back.}
\label{tab:comparison}
\end{table}

% provenance: see the header of training/train-2replies_3460844/tables_eval/sequences.tex
\begin{table}[h]
\centering\small
% GENERATED by scripts/frozen_render.py. Do not edit; re-run the renderer.
% frozen: docs/stage12_agent/stage2-4b-agent_3458956, code commit 78002066ae308d6dc6aa0ebd95f400b154fa3d49, Slurm job 3458956, model models/hf_base/iter_000, 48 tasks x 4 samples, max_tokens 1000
% ckpt1: docs/training/train-2replies_3460844/eval-train-2replies-ckpt01_3461865, code commit b80823043e95e55a41fd518734c287892743d7c4, Slurm job 3461865, model models/hf_base/iter_000 with the first 1 adapter(s) of /sc/home/nima.hamidisiboni/work/lora4vr-train-654ad1b/models/lora/train-2replies_3460844 merged in, 48 tasks x 4 samples, max_tokens 1000
% ckpt2: docs/training/train-2replies_3460844/eval-train-2replies-ckpt02_3462056, code commit b80823043e95e55a41fd518734c287892743d7c4, Slurm job 3462056, model models/hf_base/iter_000 with the first 2 adapter(s) of /sc/home/nima.hamidisiboni/work/lora4vr-train-654ad1b/models/lora/train-2replies_3460844 merged in, 48 tasks x 4 samples, max_tokens 1000
% ckpt3: docs/training/train-2replies_3460844/eval-train-2replies-ckpt03_3462280, code commit b80823043e95e55a41fd518734c287892743d7c4, Slurm job 3462280, model models/hf_base/iter_000 with the first 3 adapter(s) of /sc/home/nima.hamidisiboni/work/lora4vr-train-654ad1b/models/lora/train-2replies_3460844 merged in, 48 tasks x 4 samples, max_tokens 1000
% ckpt5: docs/training/train-2replies_3460844/eval-train-2replies-ckpt05_3470604, code commit b80823043e95e55a41fd518734c287892743d7c4, Slurm job 3470604, model models/hf_base/iter_000 with the first 5 adapter(s) of /sc/home/nima.hamidisiboni/work/lora4vr-train-654ad1b/models/lora/train-2replies_3460844 merged in, 48 tasks x 4 samples, max_tokens 1000
% ckpt23: docs/training/train-2replies_3460844/eval-train-2replies-ckpt23_3470605, code commit b80823043e95e55a41fd518734c287892743d7c4, Slurm job 3470605, model models/hf_base/iter_000 with the first 23 adapter(s) of /sc/home/nima.hamidisiboni/work/lora4vr-train-654ad1b/models/lora/train-2replies_3460844 merged in, 48 tasks x 4 samples, max_tokens 1000
\begin{tabular}{@{}lrrrrrr@{}}
\toprule
Replies of the sample, in order & frozen & ckpt1 & ckpt2 & ckpt3 & ckpt5 & ckpt23 \\
\midrule
final & 19 & 22 & 167 & 140 & 191 & 192 \\
probe, final & 38 & 140 & 14 & 36 & 1 & 0 \\
probe, probe & 132 & 23 & 7 & 16 & 0 & 0 \\
invalid call, final & 1 & 5 & 2 & 0 & 0 & 0 \\
invalid call, probe & 1 & 1 & 1 & 0 & 0 & 0 \\
invalid call, invalid call & 0 & 1 & 1 & 0 & 0 & 0 \\
probe, invalid call & 1 & 0 & 0 & 0 & 0 & 0 \\
\midrule
samples & 192 & 192 & 192 & 192 & 192 & 192 \\
\bottomrule
\end{tabular}

\caption{The samples by what their replies were, in order.}
\label{tab:sequences}
\end{table}

The held-out tasks show the same course as the training record. At
checkpoint~1 the first reply is a probe run in 163 of 192 samples, against
171 for the frozen model, and 140 samples are a probe followed by a final
answer, against 38. At checkpoint~2 the first reply is a probe run in 21
samples and the final answer in 167. Checkpoint~3 probes in 52: at the end
of iteration~3 the probe came back for a moment (two and three of four
first replies in the two steps before its last), before iteration~4 removed
it.

The reward rises from $-231.7$ to $-99.2$ and then to about $-80$, and the
whole rise is the error floor emptying: over the samples that gave a final
run, the reward is $-74.0$ for the frozen model and $-69.1$, $-65.1$ and
$-61.0$ for the three checkpoints (Table~\ref{tab:answered}). The controllers
themselves are hardly better than the frozen model's.

% provenance: see the header of training/train-2replies_3460844/tables_eval/answered.tex
\begin{table}[h]
\centering\small
% GENERATED by scripts/frozen_render.py. Do not edit; re-run the renderer.
% frozen: docs/stage12_agent/stage2-4b-agent_3458956, code commit 78002066ae308d6dc6aa0ebd95f400b154fa3d49, Slurm job 3458956, model models/hf_base/iter_000, 48 tasks x 4 samples, max_tokens 1000
% ckpt1: docs/training/train-2replies_3460844/eval-train-2replies-ckpt01_3461865, code commit b80823043e95e55a41fd518734c287892743d7c4, Slurm job 3461865, model models/hf_base/iter_000 with the first 1 adapter(s) of /sc/home/nima.hamidisiboni/work/lora4vr-train-654ad1b/models/lora/train-2replies_3460844 merged in, 48 tasks x 4 samples, max_tokens 1000
% ckpt2: docs/training/train-2replies_3460844/eval-train-2replies-ckpt02_3462056, code commit b80823043e95e55a41fd518734c287892743d7c4, Slurm job 3462056, model models/hf_base/iter_000 with the first 2 adapter(s) of /sc/home/nima.hamidisiboni/work/lora4vr-train-654ad1b/models/lora/train-2replies_3460844 merged in, 48 tasks x 4 samples, max_tokens 1000
% ckpt3: docs/training/train-2replies_3460844/eval-train-2replies-ckpt03_3462280, code commit b80823043e95e55a41fd518734c287892743d7c4, Slurm job 3462280, model models/hf_base/iter_000 with the first 3 adapter(s) of /sc/home/nima.hamidisiboni/work/lora4vr-train-654ad1b/models/lora/train-2replies_3460844 merged in, 48 tasks x 4 samples, max_tokens 1000
% ckpt5: docs/training/train-2replies_3460844/eval-train-2replies-ckpt05_3470604, code commit b80823043e95e55a41fd518734c287892743d7c4, Slurm job 3470604, model models/hf_base/iter_000 with the first 5 adapter(s) of /sc/home/nima.hamidisiboni/work/lora4vr-train-654ad1b/models/lora/train-2replies_3460844 merged in, 48 tasks x 4 samples, max_tokens 1000
% ckpt23: docs/training/train-2replies_3460844/eval-train-2replies-ckpt23_3470605, code commit b80823043e95e55a41fd518734c287892743d7c4, Slurm job 3470605, model models/hf_base/iter_000 with the first 23 adapter(s) of /sc/home/nima.hamidisiboni/work/lora4vr-train-654ad1b/models/lora/train-2replies_3460844 merged in, 48 tasks x 4 samples, max_tokens 1000
\begin{tabular}{@{}lrrrrrr@{}}
\toprule
 & & & & \multicolumn{3}{c}{The samples that gave a final run} \\
\cmidrule(l){5-7}
Run & Samples & No final run & Reward, all & Reward & Success & Fast \\
\midrule
frozen & 192 & 134 & $-$231.7 & $-$74.0 & 31.0\% & 3.4\% \\
ckpt1 & 192 & 25 & $-$99.2 & $-$69.1 & 25.7\% & 1.2\% \\
ckpt2 & 192 & 9 & $-$76.1 & $-$65.1 & 26.2\% & 2.7\% \\
ckpt3 & 192 & 16 & $-$80.9 & $-$61.0 & 24.4\% & 2.8\% \\
ckpt5 & 192 & 0 & $-$41.9 & $-$41.9 & 32.8\% & 5.2\% \\
ckpt23 & 192 & 0 & $-$41.1 & $-$41.1 & 25.5\% & 25.0\% \\
\bottomrule
\end{tabular}

\caption{Each run split into the samples that ended without a final run,
which score $-300$, and the rest.}
\label{tab:answered}
\end{table}

\section{Reading}

\begin{enumerate}[leftmargin=*, itemsep=0.5em]
\item \textbf{The first iteration did what was hoped, and more.} The samples
  that end without a final run fall from 134 to 25 of 192. The probes also
  improve: of the first-reply probe runs, 145 of 163 run without an error
  against 107 of 171, and 30 print a number within 0.01 of the hidden limit
  against 11 (Table~\ref{tab:first}). One pass over 24 tasks was enough to
  teach the model that the second reply has to be the answer.

\item \textbf{The reward said ``do not probe'' from the first step on.}
  Table~\ref{tab:value} gives the mean reward of a sample by what its first
  reply was. This is the quantity the method credits a first reply with: its
  own mean outcome less the mean over the tree. For the frozen model a first
  reply that is a probe run is worth $-249.5$ and one that is the final
  answer $-75.6$, a gap of 174 points against the probe. Learning to answer
  after the probe closed two thirds of the gap and no more: at checkpoint~1
  the probe is worth $-105.1$ and the direct answer $-51.0$. Every step of
  iteration~2 therefore pushed the first reply toward the direct answer, and
  the probe was gone within the iteration.

\item \textbf{What was left of the gap has two parts, and the second is the
  one that matters.} First, a probe can still end without an answer (23 of
  163 at checkpoint~1), and each such sample costs $-300$. Second, a sample
  that probes and then answers is not better than one that answers straight
  away: $-71.9$ against $-51.0$ at checkpoint~1 (Table~\ref{tab:path}). The
  probe does not yet buy a better controller. Where it found the limit and
  the controller carries the number, the outcome is worse than not knowing:
  of the 15 such controllers at checkpoint~1, 13 ruin the batch by a shower
  and one succeeds, and 9 of the 13 hold a number that lies above the true
  limit (Table~\ref{tab:fate}). The model reads the level at which the probe
  batch showered and rides at it, with no margin below. As in the frozen
  model, the number reaches the controller; riding safely under it is what
  is missing.

\item \textbf{So the two skills have to be learned in the wrong order.}
  For the reward to favour the probe, the model must answer after it every
  time \emph{and} turn the measured limit into a controller that beats a
  blind one. The first was learned in 24 steps. The second was not there
  when the first arrived, and from that moment the reward removed the probe,
  and with it every further sample from which the second could have been
  learned. The reply budget is part of this, since a failed probe has no
  third reply to be repaired in, but it is not all of it: with every probe
  answered, the probe would still be worth no more than the direct answer.

\item \textbf{Without the probe the model trains toward a gamble.} At
  checkpoints~2 and~3 more than half of all samples end in a ruined batch
  (113 and 117 of 192, against 30 for the frozen model, most of whose
  samples never ran; Table~\ref{tab:diag}), and the ride ratio rises from
  0.65 to 0.82 and 0.89.
  By thirds of the limit's range (Table~\ref{tab:third}), success is 59\% in
  the high third and under 2\% in the low third: the controllers cool as if
  the limit were high. The hand-written controllers that do not know the
  limit score $+38.4$ (the cautious one, which rides under the lowest limit)
  and $+46.2$ (the best one found by search in Stage~0, which gives up some
  of the low third); the run stays far below both, at $-76$ and $-81$ on the
  held-out tasks, and on the training tasks its reward falls again after
  iteration~\TRbestiteration. Whether the last checkpoints are
  worse than the fifth on the held-out tasks too is not yet measured (see
  below).

\item \textbf{What the run does not show.} It is one seed, one learning
  rate and one sizing, and the schedule restarts every iteration. It does not
  show that two replies cannot work. It shows that at the frozen model's
  starting point the reward gives no pull toward the probe, measures how
  large the push against it is, and locates the missing piece: a controller
  that keeps a margin under a measured limit.
\end{enumerate}

% provenance: see the header of training/train-2replies_3460844/tables_eval/first_reply_value.tex
\begin{table}[h]
\centering\small
% GENERATED by scripts/frozen_render.py. Do not edit; re-run the renderer.
% frozen: docs/stage12_agent/stage2-4b-agent_3458956, code commit 78002066ae308d6dc6aa0ebd95f400b154fa3d49, Slurm job 3458956, model models/hf_base/iter_000, 48 tasks x 4 samples, max_tokens 1000
% ckpt1: docs/training/train-2replies_3460844/eval-train-2replies-ckpt01_3461865, code commit b80823043e95e55a41fd518734c287892743d7c4, Slurm job 3461865, model models/hf_base/iter_000 with the first 1 adapter(s) of /sc/home/nima.hamidisiboni/work/lora4vr-train-654ad1b/models/lora/train-2replies_3460844 merged in, 48 tasks x 4 samples, max_tokens 1000
% ckpt2: docs/training/train-2replies_3460844/eval-train-2replies-ckpt02_3462056, code commit b80823043e95e55a41fd518734c287892743d7c4, Slurm job 3462056, model models/hf_base/iter_000 with the first 2 adapter(s) of /sc/home/nima.hamidisiboni/work/lora4vr-train-654ad1b/models/lora/train-2replies_3460844 merged in, 48 tasks x 4 samples, max_tokens 1000
% ckpt3: docs/training/train-2replies_3460844/eval-train-2replies-ckpt03_3462280, code commit b80823043e95e55a41fd518734c287892743d7c4, Slurm job 3462280, model models/hf_base/iter_000 with the first 3 adapter(s) of /sc/home/nima.hamidisiboni/work/lora4vr-train-654ad1b/models/lora/train-2replies_3460844 merged in, 48 tasks x 4 samples, max_tokens 1000
% ckpt5: docs/training/train-2replies_3460844/eval-train-2replies-ckpt05_3470604, code commit b80823043e95e55a41fd518734c287892743d7c4, Slurm job 3470604, model models/hf_base/iter_000 with the first 5 adapter(s) of /sc/home/nima.hamidisiboni/work/lora4vr-train-654ad1b/models/lora/train-2replies_3460844 merged in, 48 tasks x 4 samples, max_tokens 1000
% ckpt23: docs/training/train-2replies_3460844/eval-train-2replies-ckpt23_3470605, code commit b80823043e95e55a41fd518734c287892743d7c4, Slurm job 3470605, model models/hf_base/iter_000 with the first 23 adapter(s) of /sc/home/nima.hamidisiboni/work/lora4vr-train-654ad1b/models/lora/train-2replies_3460844 merged in, 48 tasks x 4 samples, max_tokens 1000
\begin{tabular}{@{}llrrrrr@{}}
\toprule
Run & The first reply is & $n$ & Reached a final run & Success & Fast & Reward \\
\midrule
frozen & a final answer & 19 & 100.0\% & 26.3\% & 5.3\% & $-$75.6 \\
 & a probe run & 171 & 22.2\% & 7.6\% & 0.6\% & $-$249.5 \\
 & a tool call that could not be acted on & 2 & 50.0\% & 0.0\% & 0.0\% & $-$193.9 \\
 & \emph{a probe run, less a final answer} & & & & & $-$173.9 \\
\addlinespace
ckpt1 & a final answer & 22 & 100.0\% & 31.8\% & 4.5\% & $-$51.0 \\
 & a probe run & 163 & 85.9\% & 20.2\% & 0.6\% & $-$105.1 \\
 & a tool call that could not be acted on & 7 & 71.4\% & 42.9\% & 0.0\% & $-$111.8 \\
 & \emph{a probe run, less a final answer} & & & & & $-$54.1 \\
\addlinespace
ckpt2 & a final answer & 167 & 100.0\% & 26.9\% & 3.0\% & $-$65.6 \\
 & a probe run & 21 & 66.7\% & 14.3\% & 0.0\% & $-$137.3 \\
 & a tool call that could not be acted on & 4 & 50.0\% & 0.0\% & 0.0\% & $-$194.7 \\
 & \emph{a probe run, less a final answer} & & & & & $-$71.7 \\
\addlinespace
ckpt3 & a final answer & 140 & 100.0\% & 25.7\% & 2.9\% & $-$60.5 \\
 & a probe run & 52 & 69.2\% & 13.5\% & 1.9\% & $-$135.8 \\
 & a tool call that could not be acted on & 0 & -- & -- & -- & -- \\
 & \emph{a probe run, less a final answer} & & & & & $-$75.4 \\
\addlinespace
ckpt5 & a final answer & 191 & 100.0\% & 33.0\% & 5.2\% & $-$41.7 \\
 & a probe run & 1 & 100.0\% & 0.0\% & 0.0\% & $-$84.8 \\
 & a tool call that could not be acted on & 0 & -- & -- & -- & -- \\
 & \emph{a probe run, less a final answer} & & & & & $-$43.1 \\
\addlinespace
ckpt23 & a final answer & 192 & 100.0\% & 25.5\% & 25.0\% & $-$41.1 \\
 & a probe run & 0 & -- & -- & -- & -- \\
 & a tool call that could not be acted on & 0 & -- & -- & -- & -- \\
 & \emph{a probe run, less a final answer} & & & & & -- \\
\bottomrule
\end{tabular}

\caption{What a first reply is worth: the mean outcome of the samples by
their first reply. The last row of each run is the reward of a probe run
less that of a final answer.}
\label{tab:value}
\end{table}

% provenance: see the header of training/train-2replies_3460844/tables_eval/by_path.tex
\begin{table}[h]
\centering\small
\resizebox{\linewidth}{!}{% GENERATED by scripts/frozen_render.py. Do not edit; re-run the renderer.
% frozen: docs/stage12_agent/stage2-4b-agent_3458956, code commit 78002066ae308d6dc6aa0ebd95f400b154fa3d49, Slurm job 3458956, model models/hf_base/iter_000, 48 tasks x 4 samples, max_tokens 1000
% ckpt1: docs/training/train-2replies_3460844/eval-train-2replies-ckpt01_3461865, code commit b80823043e95e55a41fd518734c287892743d7c4, Slurm job 3461865, model models/hf_base/iter_000 with the first 1 adapter(s) of /sc/home/nima.hamidisiboni/work/lora4vr-train-654ad1b/models/lora/train-2replies_3460844 merged in, 48 tasks x 4 samples, max_tokens 1000
% ckpt2: docs/training/train-2replies_3460844/eval-train-2replies-ckpt02_3462056, code commit b80823043e95e55a41fd518734c287892743d7c4, Slurm job 3462056, model models/hf_base/iter_000 with the first 2 adapter(s) of /sc/home/nima.hamidisiboni/work/lora4vr-train-654ad1b/models/lora/train-2replies_3460844 merged in, 48 tasks x 4 samples, max_tokens 1000
% ckpt3: docs/training/train-2replies_3460844/eval-train-2replies-ckpt03_3462280, code commit b80823043e95e55a41fd518734c287892743d7c4, Slurm job 3462280, model models/hf_base/iter_000 with the first 3 adapter(s) of /sc/home/nima.hamidisiboni/work/lora4vr-train-654ad1b/models/lora/train-2replies_3460844 merged in, 48 tasks x 4 samples, max_tokens 1000
% ckpt5: docs/training/train-2replies_3460844/eval-train-2replies-ckpt05_3470604, code commit b80823043e95e55a41fd518734c287892743d7c4, Slurm job 3470604, model models/hf_base/iter_000 with the first 5 adapter(s) of /sc/home/nima.hamidisiboni/work/lora4vr-train-654ad1b/models/lora/train-2replies_3460844 merged in, 48 tasks x 4 samples, max_tokens 1000
% ckpt23: docs/training/train-2replies_3460844/eval-train-2replies-ckpt23_3470605, code commit b80823043e95e55a41fd518734c287892743d7c4, Slurm job 3470605, model models/hf_base/iter_000 with the first 23 adapter(s) of /sc/home/nima.hamidisiboni/work/lora4vr-train-654ad1b/models/lora/train-2replies_3460844 merged in, 48 tasks x 4 samples, max_tokens 1000
\begin{tabular}{@{}llrrrrrrr@{}}
\toprule
Run & Path & $n$ & Controller ran & Success & Fast & Reward & Ride ratio & Est.\ within 0.01 \\
\midrule
frozen & answered straight away & 19 & 89.5\% & 26.3\% & 5.3\% & $-$75.6 & 0.68 & -- \\
 & probed, then a final run & 39 & 89.7\% & 33.3\% & 2.6\% & $-$73.2 & 0.64 & 6 \\
 & ended without a final run & 134 & 0.0\% & 0.0\% & 0.0\% & $-$300.0 & -- & 5 \\
\addlinespace
ckpt1 & answered straight away & 22 & 100.0\% & 31.8\% & 4.5\% & $-$51.0 & 0.59 & -- \\
 & probed, then a final run & 145 & 95.9\% & 24.8\% & 0.7\% & $-$71.9 & 0.68 & 29 \\
 & ended without a final run & 25 & 0.0\% & 0.0\% & 0.0\% & $-$300.0 & -- & 1 \\
\addlinespace
ckpt2 & answered straight away & 167 & 94.6\% & 26.9\% & 3.0\% & $-$65.6 & 0.80 & -- \\
 & probed, then a final run & 16 & 100.0\% & 18.8\% & 0.0\% & $-$60.1 & 0.95 & 1 \\
 & ended without a final run & 9 & 0.0\% & 0.0\% & 0.0\% & $-$300.0 & -- & 2 \\
\addlinespace
ckpt3 & answered straight away & 140 & 95.7\% & 25.7\% & 2.9\% & $-$60.5 & 0.88 & -- \\
 & probed, then a final run & 36 & 100.0\% & 19.4\% & 2.8\% & $-$62.9 & 0.92 & 3 \\
 & ended without a final run & 16 & 0.0\% & 0.0\% & 0.0\% & $-$300.0 & -- & 0 \\
\addlinespace
ckpt5 & answered straight away & 191 & 98.4\% & 33.0\% & 5.2\% & $-$41.7 & 0.91 & -- \\
 & probed, then a final run & 1 & 100.0\% & 0.0\% & 0.0\% & $-$84.8 & 1.00 & 1 \\
 & ended without a final run & 0 & -- & -- & -- & -- & -- & 0 \\
\addlinespace
ckpt23 & answered straight away & 192 & 100.0\% & 25.5\% & 25.0\% & $-$41.1 & 1.00 & -- \\
 & probed, then a final run & 0 & -- & -- & -- & -- & -- & 0 \\
 & ended without a final run & 0 & -- & -- & -- & -- & -- & 0 \\
\bottomrule
\end{tabular}
}
\caption{The samples by the path they took. Est.\ within 0.01: samples in
which some probe run printed a number within 0.01 of the limit.}
\label{tab:path}
\end{table}

% provenance: see the header of training/train-2replies_3460844/tables_eval/first_actions.tex
\begin{table}[h]
\centering\small
% GENERATED by scripts/frozen_render.py. Do not edit; re-run the renderer.
% frozen: docs/stage12_agent/stage2-4b-agent_3458956, code commit 78002066ae308d6dc6aa0ebd95f400b154fa3d49, Slurm job 3458956, model models/hf_base/iter_000, 48 tasks x 4 samples, max_tokens 1000
% ckpt1: docs/training/train-2replies_3460844/eval-train-2replies-ckpt01_3461865, code commit b80823043e95e55a41fd518734c287892743d7c4, Slurm job 3461865, model models/hf_base/iter_000 with the first 1 adapter(s) of /sc/home/nima.hamidisiboni/work/lora4vr-train-654ad1b/models/lora/train-2replies_3460844 merged in, 48 tasks x 4 samples, max_tokens 1000
% ckpt2: docs/training/train-2replies_3460844/eval-train-2replies-ckpt02_3462056, code commit b80823043e95e55a41fd518734c287892743d7c4, Slurm job 3462056, model models/hf_base/iter_000 with the first 2 adapter(s) of /sc/home/nima.hamidisiboni/work/lora4vr-train-654ad1b/models/lora/train-2replies_3460844 merged in, 48 tasks x 4 samples, max_tokens 1000
% ckpt3: docs/training/train-2replies_3460844/eval-train-2replies-ckpt03_3462280, code commit b80823043e95e55a41fd518734c287892743d7c4, Slurm job 3462280, model models/hf_base/iter_000 with the first 3 adapter(s) of /sc/home/nima.hamidisiboni/work/lora4vr-train-654ad1b/models/lora/train-2replies_3460844 merged in, 48 tasks x 4 samples, max_tokens 1000
% ckpt5: docs/training/train-2replies_3460844/eval-train-2replies-ckpt05_3470604, code commit b80823043e95e55a41fd518734c287892743d7c4, Slurm job 3470604, model models/hf_base/iter_000 with the first 5 adapter(s) of /sc/home/nima.hamidisiboni/work/lora4vr-train-654ad1b/models/lora/train-2replies_3460844 merged in, 48 tasks x 4 samples, max_tokens 1000
% ckpt23: docs/training/train-2replies_3460844/eval-train-2replies-ckpt23_3470605, code commit b80823043e95e55a41fd518734c287892743d7c4, Slurm job 3470605, model models/hf_base/iter_000 with the first 23 adapter(s) of /sc/home/nima.hamidisiboni/work/lora4vr-train-654ad1b/models/lora/train-2replies_3460844 merged in, 48 tasks x 4 samples, max_tokens 1000
\begin{tabular}{@{}lrrrrrr@{}}
\toprule
 & frozen & ckpt1 & ckpt2 & ckpt3 & ckpt5 & ckpt23 \\
\midrule
samples & 192 & 192 & 192 & 192 & 192 & 192 \\
\midrule
\multicolumn{7}{@{}l}{the first reply is} \\
\quad a final run (the controller straight away) & 19 & 22 & 167 & 140 & 191 & 192 \\
\quad a probe run, whose output came back & 171 & 163 & 21 & 52 & 1 & 0 \\
\quad \quad the probe code ran without an error & 107 & 145 & 18 & 38 & 1 & 0 \\
\quad \quad the output holds a number within 0.01 of the limit & 11 & 30 & 3 & 3 & 1 & 0 \\
\quad \quad the task went on to a final run & 38 & 140 & 14 & 36 & 1 & 0 \\
\quad a tool call that could not be acted on & 2 & 7 & 4 & 0 & 0 & 0 \\
\quad \quad cut off by the token budget & 2 & 6 & 4 & 0 & 0 & 0 \\
\quad \quad not valid JSON & 0 & 1 & 0 & 0 & 0 & 0 \\
\quad \quad other (unknown tool, no code) & 0 & 0 & 0 & 0 & 0 & 0 \\
\midrule
samples that ended without a final run & 134 & 25 & 9 & 16 & 0 & 0 \\
samples in which some reply mentions the reply budget & 0 & 0 & 0 & 0 & 0 & 0 \\
\bottomrule
\end{tabular}

\caption{The first reply of every sample.}
\label{tab:first}
\end{table}

% provenance: see the header of training/train-2replies_3460844/tables_eval/diagnostics.tex
\begin{table}[h]
\centering\small
\resizebox{\linewidth}{!}{% GENERATED by scripts/frozen_render.py. Do not edit; re-run the renderer.
% frozen: docs/stage12_agent/stage2-4b-agent_3458956, code commit 78002066ae308d6dc6aa0ebd95f400b154fa3d49, Slurm job 3458956, model models/hf_base/iter_000, 48 tasks x 4 samples, max_tokens 1000
% ckpt1: docs/training/train-2replies_3460844/eval-train-2replies-ckpt01_3461865, code commit b80823043e95e55a41fd518734c287892743d7c4, Slurm job 3461865, model models/hf_base/iter_000 with the first 1 adapter(s) of /sc/home/nima.hamidisiboni/work/lora4vr-train-654ad1b/models/lora/train-2replies_3460844 merged in, 48 tasks x 4 samples, max_tokens 1000
% ckpt2: docs/training/train-2replies_3460844/eval-train-2replies-ckpt02_3462056, code commit b80823043e95e55a41fd518734c287892743d7c4, Slurm job 3462056, model models/hf_base/iter_000 with the first 2 adapter(s) of /sc/home/nima.hamidisiboni/work/lora4vr-train-654ad1b/models/lora/train-2replies_3460844 merged in, 48 tasks x 4 samples, max_tokens 1000
% ckpt3: docs/training/train-2replies_3460844/eval-train-2replies-ckpt03_3462280, code commit b80823043e95e55a41fd518734c287892743d7c4, Slurm job 3462280, model models/hf_base/iter_000 with the first 3 adapter(s) of /sc/home/nima.hamidisiboni/work/lora4vr-train-654ad1b/models/lora/train-2replies_3460844 merged in, 48 tasks x 4 samples, max_tokens 1000
\begin{tabular}{@{}lrrrrrrrrrrr@{}}
\toprule
Run & $n$ & Reply chars & Cut off & Success & Ruined & Timeout & Error & Probe chars & Probe crashed & Probe cut off & Est.\ within 0.01 \\
\midrule
frozen & 192 & 1496 & 4.2\% & 18 & 30 & 4 & 140 & 529--3803 & 36.3\% & 0.0\% & 6.4\% \\
ckpt1 & 192 & 2303 & 3.1\% & 43 & 103 & 15 & 31 & 833--3450 & 8.0\% & 0.0\% & 18.4\% \\
ckpt2 & 192 & 2774 & 10.4\% & 48 & 113 & 13 & 18 & 1436--3891 & 14.3\% & 0.0\% & 14.3\% \\
ckpt3 & 192 & 2196 & 4.2\% & 43 & 117 & 10 & 22 & 805--3100 & 26.9\% & 0.0\% & 5.8\% \\
\bottomrule
\end{tabular}
}
\caption{How the samples ended (success, ruined, timeout, error: counts of
192) and what the replies were like. Reply chars: the median length of the
last reply. Probe crashed and Est.\ within 0.01: of the samples whose first
reply was a probe run, the share whose tool response reports an error and
the share whose output holds a number within 0.01 of the limit.}
\label{tab:diag}
\end{table}

% provenance: see the header of training/train-2replies_3460844/tables_eval/by_third.tex
\begin{table}[h]
\centering\small
\resizebox{\linewidth}{!}{% GENERATED by scripts/frozen_render.py. Do not edit; re-run the renderer.
% frozen: docs/stage12_agent/stage2-4b-agent_3458956, code commit 78002066ae308d6dc6aa0ebd95f400b154fa3d49, Slurm job 3458956, model models/hf_base/iter_000, 48 tasks x 4 samples, max_tokens 1000
% ckpt1: docs/training/train-2replies_3460844/eval-train-2replies-ckpt01_3461865, code commit b80823043e95e55a41fd518734c287892743d7c4, Slurm job 3461865, model models/hf_base/iter_000 with the first 1 adapter(s) of /sc/home/nima.hamidisiboni/work/lora4vr-train-654ad1b/models/lora/train-2replies_3460844 merged in, 48 tasks x 4 samples, max_tokens 1000
% ckpt2: docs/training/train-2replies_3460844/eval-train-2replies-ckpt02_3462056, code commit b80823043e95e55a41fd518734c287892743d7c4, Slurm job 3462056, model models/hf_base/iter_000 with the first 2 adapter(s) of /sc/home/nima.hamidisiboni/work/lora4vr-train-654ad1b/models/lora/train-2replies_3460844 merged in, 48 tasks x 4 samples, max_tokens 1000
% ckpt3: docs/training/train-2replies_3460844/eval-train-2replies-ckpt03_3462280, code commit b80823043e95e55a41fd518734c287892743d7c4, Slurm job 3462280, model models/hf_base/iter_000 with the first 3 adapter(s) of /sc/home/nima.hamidisiboni/work/lora4vr-train-654ad1b/models/lora/train-2replies_3460844 merged in, 48 tasks x 4 samples, max_tokens 1000
% ckpt5: docs/training/train-2replies_3460844/eval-train-2replies-ckpt05_3470604, code commit b80823043e95e55a41fd518734c287892743d7c4, Slurm job 3470604, model models/hf_base/iter_000 with the first 5 adapter(s) of /sc/home/nima.hamidisiboni/work/lora4vr-train-654ad1b/models/lora/train-2replies_3460844 merged in, 48 tasks x 4 samples, max_tokens 1000
% ckpt23: docs/training/train-2replies_3460844/eval-train-2replies-ckpt23_3470605, code commit b80823043e95e55a41fd518734c287892743d7c4, Slurm job 3470605, model models/hf_base/iter_000 with the first 23 adapter(s) of /sc/home/nima.hamidisiboni/work/lora4vr-train-654ad1b/models/lora/train-2replies_3460844 merged in, 48 tasks x 4 samples, max_tokens 1000
\begin{tabular}{@{}llrrrrrrrrrrrrrrrrrr@{}}
\toprule
 & & \multicolumn{3}{c}{frozen} & \multicolumn{3}{c}{ckpt1} & \multicolumn{3}{c}{ckpt2} & \multicolumn{3}{c}{ckpt3} & \multicolumn{3}{c}{ckpt5} & \multicolumn{3}{c}{ckpt23} \\
\cmidrule(lr){3-5}\cmidrule(lr){6-8}\cmidrule(lr){9-11}\cmidrule(lr){12-14}\cmidrule(lr){15-17}\cmidrule(lr){18-20}
Third & $n$ & Fast & Succ. & Reward & Fast & Succ. & Reward & Fast & Succ. & Reward & Fast & Succ. & Reward & Fast & Succ. & Reward & Fast & Succ. & Reward \\
\midrule
low & 64 & 0.0\% & 3.1\% & $-$235.3 & 0.0\% & 6.2\% & $-$121.4 & 0.0\% & 1.6\% & $-$104.0 & 0.0\% & 1.6\% & $-$101.7 & 0.0\% & 0.0\% & $-$87.2 & 0.0\% & 0.0\% & $-$85.4 \\
middle & 64 & 0.0\% & 6.2\% & $-$243.9 & 0.0\% & 23.4\% & $-$85.2 & 1.6\% & 14.1\% & $-$96.5 & 0.0\% & 6.2\% & $-$103.4 & 3.1\% & 10.9\% & $-$81.5 & 0.0\% & 0.0\% & $-$84.8 \\
high & 64 & 3.1\% & 18.8\% & $-$216.0 & 3.1\% & 37.5\% & $-$91.0 & 6.2\% & 59.4\% & $-$27.9 & 7.8\% & 59.4\% & $-$37.5 & 12.5\% & 87.5\% & 42.8 & 75.0\% & 76.6\% & 47.0 \\
\bottomrule
\end{tabular}
}
\caption{Outcome by thirds of the range of the hidden limit (low: 0.08 to
0.137; high: 0.193 to 0.25).}
\label{tab:third}
\end{table}

% provenance: see the header of training/train-2replies_3460844/tables_eval/probe_fate.tex
\begin{table}[h]
\centering\small
\resizebox{0.92\linewidth}{!}{% GENERATED by scripts/frozen_render.py. Do not edit; re-run the renderer.
% frozen: docs/stage12_agent/stage2-4b-agent_3458956, code commit 78002066ae308d6dc6aa0ebd95f400b154fa3d49, Slurm job 3458956, model models/hf_base/iter_000, 48 tasks x 4 samples, max_tokens 1000
% ckpt1: docs/training/train-2replies_3460844/eval-train-2replies-ckpt01_3461865, code commit b80823043e95e55a41fd518734c287892743d7c4, Slurm job 3461865, model models/hf_base/iter_000 with the first 1 adapter(s) of /sc/home/nima.hamidisiboni/work/lora4vr-train-654ad1b/models/lora/train-2replies_3460844 merged in, 48 tasks x 4 samples, max_tokens 1000
% ckpt2: docs/training/train-2replies_3460844/eval-train-2replies-ckpt02_3462056, code commit b80823043e95e55a41fd518734c287892743d7c4, Slurm job 3462056, model models/hf_base/iter_000 with the first 2 adapter(s) of /sc/home/nima.hamidisiboni/work/lora4vr-train-654ad1b/models/lora/train-2replies_3460844 merged in, 48 tasks x 4 samples, max_tokens 1000
% ckpt3: docs/training/train-2replies_3460844/eval-train-2replies-ckpt03_3462280, code commit b80823043e95e55a41fd518734c287892743d7c4, Slurm job 3462280, model models/hf_base/iter_000 with the first 3 adapter(s) of /sc/home/nima.hamidisiboni/work/lora4vr-train-654ad1b/models/lora/train-2replies_3460844 merged in, 48 tasks x 4 samples, max_tokens 1000
% ckpt5: docs/training/train-2replies_3460844/eval-train-2replies-ckpt05_3470604, code commit b80823043e95e55a41fd518734c287892743d7c4, Slurm job 3470604, model models/hf_base/iter_000 with the first 5 adapter(s) of /sc/home/nima.hamidisiboni/work/lora4vr-train-654ad1b/models/lora/train-2replies_3460844 merged in, 48 tasks x 4 samples, max_tokens 1000
% ckpt23: docs/training/train-2replies_3460844/eval-train-2replies-ckpt23_3470605, code commit b80823043e95e55a41fd518734c287892743d7c4, Slurm job 3470605, model models/hf_base/iter_000 with the first 23 adapter(s) of /sc/home/nima.hamidisiboni/work/lora4vr-train-654ad1b/models/lora/train-2replies_3460844 merged in, 48 tasks x 4 samples, max_tokens 1000
\begin{tabular}{@{}lrrrrrr@{}}
\toprule
 & frozen & ckpt1 & ckpt2 & ckpt3 & ckpt5 & ckpt23 \\
\midrule
probe runs executed & 171 & 163 & 21 & 52 & 1 & 0 \\
\quad whose output holds a number within 0.01 of the limit & 11 & 30 & 3 & 3 & 1 & 0 \\
\qquad the code ran to its end & 11 & 30 & 3 & 3 & 1 & 0 \\
\qquad the code stopped on an error after printing it & 0 & 0 & 0 & 0 & 0 & 0 \\
\midrule
samples in which a probe run printed such a number & 11 & 30 & 3 & 3 & 1 & 0 \\
\quad ran out of replies before a final answer & 5 & 1 & 2 & 0 & 0 & 0 \\
\quad went on to a final run & 6 & 29 & 1 & 3 & 1 & 0 \\
\qquad the controller holds a number within 0.01 of the limit & 4 & 15 & 0 & 2 & 0 & 0 \\
\qquad\quad the final run succeeded & 1 & 1 & 0 & 0 & 0 & 0 \\
\qquad\quad \quad and was fast & 0 & 0 & 0 & 0 & 0 & 0 \\
\qquad\quad the batch was ruined & 2 & 13 & 0 & 2 & 0 & 0 \\
\qquad\quad \quad by a shower: the reading went over the limit & 1 & 13 & 0 & 2 & 0 & 0 \\
\qquad\quad \qquad the number in the controller is above the limit & 1 & 9 & 0 & 1 & 0 & 0 \\
\qquad\quad timed out or did not run & 1 & 1 & 0 & 0 & 0 & 0 \\
\qquad the controller holds no such number & 2 & 14 & 1 & 1 & 1 & 0 \\
\qquad\quad the final run succeeded & 1 & 3 & 0 & 0 & 0 & 0 \\
\qquad\quad \quad and was fast & 0 & 0 & 0 & 0 & 0 & 0 \\
\qquad\quad the batch was ruined & 1 & 10 & 1 & 1 & 1 & 0 \\
\qquad\quad \quad by a shower: the reading went over the limit & 1 & 10 & 1 & 1 & 1 & 0 \\
\qquad\quad timed out or did not run & 0 & 1 & 0 & 0 & 0 & 0 \\
\midrule
samples that probed, found no such number, and gave a final run & 32 & 111 & 13 & 33 & 0 & 0 \\
\quad the controller holds a number within 0.01 of the limit all the same & 9 & 14 & 3 & 6 & 0 & 0 \\
\quad the final run succeeded & 11 & 29 & 3 & 7 & 0 & 0 \\
\quad \quad and was fast & 1 & 1 & 0 & 1 & 0 & 0 \\
\quad the batch was ruined & 14 & 66 & 10 & 23 & 0 & 0 \\
\quad \quad by a shower: the reading went over the limit & 9 & 40 & 10 & 23 & 0 & 0 \\
\quad timed out or did not run & 7 & 16 & 0 & 3 & 0 & 0 \\
\bottomrule
\end{tabular}
}
\caption{The probes that printed a number within 0.01 of the limit, followed
to the controller and the final run.}
\label{tab:fate}
\end{table}

\clearpage
\section{Open}

\begin{enumerate}[leftmargin=*, itemsep=0.3em]
\item \textbf{Checkpoints 5 and 23.} The best iteration by the training
  record and the last. Their evaluations are submitted and wait for the A30
  node, which is in maintenance until 2026-11-28. They say whether the fall
  of the reward after iteration~\TRbestiteration{} shows on the held-out
  tasks, and how far a run without the probe gets.
\item \textbf{Is this run, in effect, the one-reply stage?} From iteration~5
  on the model answers with its first reply, as it would with one reply and
  no tool. Stage~3 of the plan would say whether training with one reply
  reaches the same level and the same gamble.
\item \textbf{How to keep the probe until it pays, with the reward as the
  only signal and the brief untouched.} This is the question the run leaves.
  A longer reply budget removes the first part of the gap (a probe that goes
  wrong can be repaired and still answered) and not the second. The second
  part is closed only when some controllers written after a probe ride
  safely under the measured limit and are seen to beat the blind ones; how
  often the frozen model produces such a controller, and how many
  conversations per step it would take to see one, can be counted from the
  Stage~2 samples before any further run.
\end{enumerate}

\section{Where things are}

\begin{sloppypar}
\begin{itemize}[leftmargin=*, itemsep=0.2em]
\item The adapters, all \TRiterations: on the cluster, in the run's code
  directory,
  \nolinkurl{~/work/lora4vr-train-654ad1b/models/lora/train-2replies_3460844/},
  as \texttt{iter\_000} to \texttt{iter\_022} (47~MB each, 1.3~GB with the
  tokenizer copies). They are not in the repository. Checkpoint $k$ is the
  base model \texttt{models/hf\_base/iter\_000} with \texttt{iter\_000} to
  \texttt{iter\_$(k{-}1)$} merged in order.
\item The record of the run:
  \nolinkurl{docs/training/train-2replies_3460844/train/}: the config, one
  file of per-step entries per iteration, the log, and
  \texttt{completions.tgz}, every conversation of the run with the credit
  it was given, as the trainer saved it (one Parquet file per step).
\item The evaluations:
  \nolinkurl{docs/training/train-2replies_3460844/eval-train-2replies-ckpt<kk>_<job>/},
  each with the job log, the summary and every sample with its transcript.
\end{itemize}
\end{sloppypar}

\section{Reproduce}
\sloppy
% commits and jobs: see the provenance block at the top of this file
A training run gets its own code directory on the cluster: the loop starts a
new process every iteration, so files pushed under a running job would
change the run.
\begin{scriptsize}
\begin{verbatim}
C=$(git rev-parse HEAD); S=$(git rev-parse --short HEAD); D=lora4vr-train-$S
# the timing run (four steps) and the training run
hpc push . lora4vr
hpc submit lora4vr jobs/hpi/train_timing.sbatch --job-name=train-timing-a30 \
    --export=ALL,LORA4VR_GIT_COMMIT=$C
hpc push . $D
hpc sh "cd ~/work/$D && mkdir -p models logs && \
    ln -sfn ~/work/lora4vr/models/hf_base models/hf_base && uv sync --frozen"
hpc submit $D jobs/hpi/train.sbatch --job-name=train-2replies \
    --export=ALL,LORA4VR_GIT_COMMIT=$C
# one checkpoint (k iterations) on the held-out tasks, on another GPU
L=/sc/home/<user>/work/$D/models/lora/train-2replies_<jobid>
hpc submit lora4vr jobs/hpi/eval_checkpoint.sbatch --job-name=eval-train-2replies-ckpt01 \
    --export=ALL,LORA4VR_GIT_COMMIT=$C,N_ADAPTERS=1,LORA_DIR=$L
# results and tables
R=docs/training/train-2replies_<jobid>
hpc pull $D/results/train-2replies_<jobid> $R        # then: mv $R/train-2replies_<jobid> $R/train
hpc pull lora4vr/results/eval-train-2replies-ckpt01_<jobid> $R
hpc pull lora4vr/results/train-timing-a30_<jobid> docs/train_timing
uv run python scripts/train_render.py --run $R/train \
    --timing docs/train_timing/train-timing-a30_<jobid> --out $R/tables
uv run python scripts/frozen_render.py --out $R/tables_eval --macro-prefix TRE \
    frozen=docs/stage12_agent/stage2-4b-agent_<jobid> \
    ckpt1=$R/eval-train-2replies-ckpt01_<jobid> ckpt2=$R/eval-train-2replies-ckpt02_<jobid> \
    ckpt3=$R/eval-train-2replies-ckpt03_<jobid>
\end{verbatim}
\end{scriptsize}
The log of a training job holds a table of every conversation after each
step; the copy in the repository is the log without those tables, and the
conversations themselves are in \texttt{completions.tgz}.
\fussy

\clearpage
\appendix
\section{The training record, every iteration}

% provenance: see the header of training/train-2replies_3460844/tables/record.tex
\begin{table}[h]
\centering\small
\resizebox{\linewidth}{!}{% GENERATED by scripts/train_render.py. Do not edit; re-run the renderer.
% run: docs/training/train-2replies_3460844/train, code commit 654ad1bc641fbd42a46f392fd335fde2e8c134f6, Slurm job 3460844, config loop_grpo_branched_crystallizer.yaml, 23 iterations, 552 steps
% timing run: docs/train_timing/train-timing-a30_3460791, code commit 50df6169756a9047299450d1b305c54a470715a1, Slurm job 3460791, 4 steps
\begin{tabular}{@{}rrrrrrr@{}}
\toprule
Iteration & Steps & First replies followed by a second & At the error floor & Reward & Seconds per step & Peak memory (GiB) \\
\midrule
1 & 24 & 86 of 96 & 45\% & $-$166.0 & 183 & 17.4 \\
2 & 24 & 37 of 96 & 6\% & $-$62.4 & 132 & 17.9 \\
3 & 24 & 17 of 96 & 9\% & $-$66.9 & 89 & 16.7 \\
4 & 24 & 4 of 96 & 2\% & $-$48.4 & 58 & 15.5 \\
5 & 24 & 1 of 96 & 1\% & $-$26.6 & 54 & 16.1 \\
6 & 24 & 3 of 96 & 2\% & $-$37.3 & 54 & 14.3 \\
7 & 24 & 1 of 96 & 0\% & $-$37.2 & 52 & 16.6 \\
8 & 24 & 0 of 96 & 0\% & $-$36.6 & 47 & 14.8 \\
9 & 24 & 0 of 96 & 0\% & $-$49.0 & 46 & 15.3 \\
10 & 24 & 0 of 96 & 0\% & $-$45.9 & 47 & 14.3 \\
11 & 24 & 4 of 96 & 1\% & $-$39.7 & 61 & 16.6 \\
12 & 24 & 1 of 96 & 2\% & $-$40.1 & 53 & 16.0 \\
13 & 24 & 0 of 96 & 0\% & $-$46.3 & 51 & 15.8 \\
14 & 24 & 0 of 96 & 0\% & $-$51.0 & 44 & 15.3 \\
15 & 24 & 0 of 96 & 0\% & $-$62.1 & 45 & 15.0 \\
16 & 24 & 0 of 96 & 0\% & $-$56.8 & 41 & 13.4 \\
17 & 24 & 0 of 96 & 0\% & $-$55.1 & 41 & 13.8 \\
18 & 24 & 0 of 96 & 0\% & $-$54.4 & 47 & 16.3 \\
19 & 24 & 0 of 96 & 0\% & $-$56.1 & 43 & 14.4 \\
20 & 24 & 3 of 96 & 0\% & $-$51.0 & 50 & 16.4 \\
21 & 24 & 2 of 96 & 0\% & $-$57.1 & 45 & 13.8 \\
22 & 24 & 0 of 96 & 0\% & $-$55.9 & 41 & 13.4 \\
23 & 24 & 0 of 96 & 0\% & $-$54.3 & 41 & 13.5 \\
\midrule
all & 552 & 159 of 2208 & 3\% & $-$54.6 & 59 & 17.9 \\
\bottomrule
\end{tabular}
}
\caption{One row per iteration. First replies followed by a second: probe
runs, and the rare tool call that could not be acted on. At the error floor:
the share of the step's \TRconversations{} conversations without a final run
that ran, averaged over the iteration. Reward: the mean over the step's
conversations, averaged over the iteration.}
\label{tab:full}
\end{table}

\end{document}
